% !TITLE 商山早行
% !AUTHOR 温庭筠
% !DYNASTY 唐
% !GENRE 五言律诗
% !ORDER 52

\begin{poem}
晨起动征铎\textsuperscript{1}，客行悲故乡。
鸡声茅店月\textsuperscript{2}，人迹板桥霜\textsuperscript{3}。
槲叶落山路\textsuperscript{4}，枳花明驿墙\textsuperscript{5}。
因思杜陵梦\textsuperscript{6}，凫雁满回塘\textsuperscript{7}。
\end{poem}

\begin{annotations}
\notetext{1}{征铎：远行马匹系在颈上的铃铛。铎（duó），大铃。清晨起身，铃声一响，又要上路了。}
\notetext{2}{鸡声茅店月：鸡鸣声里，残月还挂在茅店上空。纯用名词，构成一幅早行图。}
\notetext{3}{人迹板桥霜：旅人的脚印印在木板桥的寒霜上——天还未亮，霜上已经有了第一行足迹。}
\notetext{4}{槲叶落山路：槲（hú）树的叶子飘落在山路上。槲叶秋天枯而不落，到春末新叶长出时才掉。}
\notetext{5}{枳花明驿墙：枳（zhǐ）树的白花在驿站墙边明亮耀眼。"明"字写花色在昏暗晨光中格外醒目。}
\notetext{6}{杜陵：长安附近的地名，诗人的故乡所在。昨夜梦回杜陵——以梦中的温暖反衬眼前的寒冷。}
\notetext{7}{凫雁满回塘：野鸭和大雁游弋在曲折的池塘中。凫（fú），野鸭。梦中的故乡温暖热闹，醒来却是孤身一人。}
\end{annotations}

\begin{appreciation}
这是一首写早行的诗，也是一首写漂泊的诗。温庭筠一生困顿，久困科场，辗转幕府，对羁旅之苦体会尤深。

铎（duó）是系在马颈上的铃铛，征铎即远行之铃。清晨起身，铃声一动，便知道又要上路了。客行悲故乡——五个字直截了当，点出全诗的底色。这悲不是嚎啕，是沉默的、习惯性的，像一件穿旧了的衣裳。

鸡声茅店月，人迹板桥霜——这是唐诗中最著名的早行名句，也是全诗最精警的一联。鸡声、茅店、月，是听觉与视觉的交织；人迹、板桥、霜，是足迹与寒意的重叠。六个名词并置，没有一个动词，没有一个形容词，却构成了一幅完整的画面：天还未亮，残月挂在茅店上空，鸡鸣声里，旅人踩着板桥上的寒霜，留下了第一行脚印。这种纯用名词的写法，后来被称为"意象叠加"，像电影里的蒙太奇，省去了一切过渡，直接把最本质的瞬间推到读者眼前。

槲（hú）叶落山路，枳（zhǐ）花明驿（yì）墙——槲树的叶子落在山路上，枳树的白花在驿站的墙边发亮。明字用得极好，在昏暗的晨光中，那一点白色格外醒目，像黑暗中擦亮的一根火柴。落叶是秋的萧瑟，明花是早的冷清，一切都在暗示：这是一段孤独的旅程。

因思杜陵梦，凫（fú）雁满回塘——于是想起昨夜梦回杜陵，凫雁成群，在曲折的池塘里游弋。杜陵是长安附近的地名，也是诗人的故乡所在。梦中的故乡温暖、丰足、热闹，而眼前的现实是寒霜、落叶、孤旅。以暖梦衬寒行，以团圆衬漂泊，这中间的落差，便是全诗最深的叹息。

温庭筠没有写自己有多想家，他只是把早行的场景一一铺陈，让读者自己去感受那份寒意。最动人的思念，往往不需要直说。
\end{appreciation}
